\section{Task diversity decides what crystallises first}
\label{sec:agenda}

\begin{figure}[t]
  \centering
  \includegraphics[width=0.66\textwidth]{agenda_trajectories}
  \caption{Trajectories in the $(\Brho, \Bet)$ plane for nine task diversities
    $\alpha = P/K$, from a random initialisation with roughly half of all triples
    frustrated. Every trajectory ends near the balanced corner: learning anneals
    frustration away. \emph{Which side of the diagonal it takes to get there is
    set by $\alpha$.}
    \todo{axis labels still read $B_I$/$B_A$; regenerate.}}
  \label{fig:trajectories}
\end{figure}

This section is independent of the spurious feature --- it holds at $d = 0$ --- and
is, we think, the most transferable result in the paper.

Start the population from random weights and random trust, so that roughly half of
all triples are unbalanced and $\Brho \approx \Bet \approx 0$. Learning drives
both balance measures toward $1$: by moving out of the dissonant quadrants of
Eq.~\ref{eq:evidence}, agents anneal frustration away. The endpoint is not the
interesting part. The \emph{path} is.

\paragraph{The result.} For small $\alpha = P/K$ --- few issues relative to the
representation dimension --- the trajectory runs \emph{above} the diagonal:
$\Bet$ rises first, and the opinion sector follows. For large $\alpha$ the
trajectory runs \emph{below} it: opinions organise first and trust follows. The
crossing is monotone in $\alpha$ and the reversal happens near $\alpha \sim 1$.

\paragraph{What it means.} A population of learners exchanging supervision
develops two kinds of structure --- agreement about the tasks, and a relational
structure over sources --- and which of them forms first is not a property of the
learning rule but of the task distribution. When the agenda is narrow, there is
not enough task signal to organise representations, so the population organises
its trust relations first and those relations then pull representations into
alignment. When the agenda is broad, representations have enough to lock onto and
trust is inferred from the resulting agreement. The symmetry
$F_w(x,y) = F_\mu(y,x)$ noted in \S\ref{sec:rule} is why both orders are available
to the same dynamics; $\alpha$ selects between them.

\paragraph{Why it matters for the rest of the paper.} It explains the difference
between the two rows of Fig.~\ref{fig:correlations}. At low $\alpha$ and strong
$d$, $\Rcw$ is small while $\Rmc$ saturates --- region (IV) --- because the trust
structure formed first and never needed the opinion sector to agree with it. The
two discriminatory regimes of \S\ref{sec:phase} are therefore not two arbitrary
patches of a diagram; they are the same two orders of crystallisation, selected in
one case by $\alpha$ and in the other by $d$.

\paragraph{Practical reading.} A system whose agents exchange supervision over a
narrow task distribution will form relational structure before it forms agreement,
and that structure will be available for a spurious feature to align itself with
before there is enough task signal to contradict it. Narrow agendas are the
dangerous regime, and they are the common one.

\todo{state the $\alpha$ values and the measurement time in-text. The source draft
leaves $\Delta t$ as a literal placeholder and it is load-bearing here: because
the dynamics anneals rather than equilibrating (\S\ref{sec:rule}), a trajectory
figure is meaningless without the time at which it was cut.}
